\chapter{MLOps + DataOps for E-Commerce}

\section{Learning objectives}
\begin{itemize}
  \item Explain why production ML in e-commerce requires disciplined operational practices beyond model training.
  \item Distinguish MLOps concerns from DataOps concerns while understanding their overlap in end-to-end commerce systems.
  \item Analyze training-serving skew, feature consistency, monitoring, drift, and deployment risk in production ML workflows.
  \item Design an ML system artifact set that includes reproducibility, observability, and governance requirements.
\end{itemize}

\section{Why MLOps and DataOps matter in e-commerce}
E-commerce organizations often move quickly from a notebook proof-of-concept to a production dependency.
Models may begin as experiments for recommendation, fraud, pricing, or service automation, but once deployed they affect revenue, customer experience, and operational load.

At that point, the technical question is no longer ``does the model score well offline?''
It becomes:
\begin{itemize}
  \item How is the model trained repeatedly?
  \item How are features computed consistently?
  \item How is the model deployed safely?
  \item How is degradation detected?
  \item How is the system rolled back when behavior becomes harmful?
\end{itemize}

MLOps addresses the operational lifecycle of models.
DataOps addresses the operational lifecycle of the data systems those models depend on.
In e-commerce, the two are tightly coupled because behavior, catalog state, pricing, inventory, and fraud patterns all change continuously.

\section{From experimentation to production systems}
Many teams can build a model.
Far fewer can operate one reliably.

\subsection{The production gap}
The production gap often appears when a promising model meets real constraints such as:
\begin{itemize}
  \item missing or delayed features,
  \item inconsistent identifiers across systems,
  \item changes in business process that alter label meaning,
  \item latency requirements,
  \item lack of rollback or approval workflow,
  \item weak monitoring once the model is live.
\end{itemize}

\subsection{Why e-commerce is especially operational}
E-commerce ML systems are unusually exposed to changing conditions:
\begin{itemize}
  \item product assortments change,
  \item promotions alter behavior,
  \item demand shifts seasonally,
  \item fraudsters adapt,
  \item support policies evolve,
  \item customers interact across multiple channels.
\end{itemize}

This means that a model that looked strong last month may become weak or unsafe if the surrounding system changes.

\section{What MLOps covers}
MLOps is the discipline of making ML systems buildable, deployable, observable, and governable.

\subsection{Typical MLOps responsibilities}
\begin{itemize}
  \item versioning code, data references, and model artifacts,
  \item training pipelines,
  \item model packaging and deployment,
  \item experiment tracking,
  \item model registry and release control,
  \item performance monitoring,
  \item rollback and incident response.
\end{itemize}

\subsection{What DataOps covers}
DataOps focuses more on:
\begin{itemize}
  \item data pipeline reliability,
  \item data contracts and validation,
  \item data freshness and lineage,
  \item reproducible transformations,
  \item quality checks across analytical and operational data movement.
\end{itemize}

In practice, production ML depends on both.
Poor DataOps can silently break a model.
Poor MLOps can deploy a model without control or visibility.

\section{Reproducibility}
Reproducibility is a foundation of trustworthy ML operations.
If the team cannot explain how a model was trained and with which data, governance and debugging become weak.

\subsection{What should be reproducible}
A strong production workflow should make it possible to recover:
\begin{itemize}
  \item source code version,
  \item feature definitions,
  \item training data snapshot or reference,
  \item hyperparameters,
  \item evaluation outputs,
  \item model artifact,
  \item deployment configuration.
\end{itemize}

\subsection{Why this matters in commerce}
When a recommender starts pushing low-margin items, or a fraud model begins over-blocking legitimate customers, teams need to determine whether the cause was:
\begin{itemize}
  \item new data,
  \item a model change,
  \item a feature transformation bug,
  \item a serving mismatch,
  \item or a business-rule change upstream.
\end{itemize}

Without reproducibility, that diagnosis becomes slow and politically difficult.

\section{Pipelines and orchestration}
Production ML systems depend on pipelines rather than ad hoc manual execution.

\subsection{Typical pipeline stages}
A simplified pipeline may include:
\begin{enumerate}
  \item ingest and validate raw data,
  \item transform features,
  \item train models,
  \item evaluate and compare candidates,
  \item register approved artifacts,
  \item deploy to staging or production,
  \item monitor outcomes and trigger retraining or rollback when needed.
\end{enumerate}

\subsection{Batch and streaming coexistence}
E-commerce ML often requires both:
\begin{itemize}
  \item batch pipelines for training and aggregate features,
  \item streaming or near-real-time pipelines for serving-time context such as session behavior, current inventory, or recent fraud indicators.
\end{itemize}

The system design must make clear which features belong to each path.

\section{Feature consistency and training-serving skew}
One of the most important production risks in ML systems is training-serving skew.
This occurs when the features or transformations used at serving time do not match those used in training.

\subsection{How skew happens}
Typical causes include:
\begin{itemize}
  \item using different code paths offline and online,
  \item stale or delayed online feature values,
  \item inconsistent joins or identifiers,
  \item business-rule changes applied only in one environment,
  \item time-window leakage in offline training that is impossible online.
\end{itemize}

\subsection{Why skew is dangerous}
Skew can degrade a model even when the model artifact itself is correct.
For example:
\begin{itemize}
  \item a fraud model may receive delayed account-age features online,
  \item a recommender may use a stale inventory flag,
  \item a pricing model may be trained on clean margin data but served with inconsistent cost feeds.
\end{itemize}

\subsection{Feature stores as a concept}
Feature stores are often introduced to reduce this risk by standardizing feature definitions and access patterns across training and serving.
The technology choice matters less than the architectural goal: one governed definition of each important feature.

\section{Model registry and release management}
Production ML systems need a controlled way to manage candidate and approved models.

\subsection{Why a registry matters}
A model registry or equivalent control plane helps track:
\begin{itemize}
  \item which model versions exist,
  \item which one is approved,
  \item which environment each version is deployed in,
  \item what evaluation evidence supported release,
  \item how to roll back if needed.
\end{itemize}

\subsection{Release strategies}
Common release strategies include:
\begin{itemize}
  \item shadow deployment,
  \item canary rollout,
  \item staged traffic ramp-up,
  \item champion-challenger comparison,
  \item full rollback on guardrail breach.
\end{itemize}

These practices are especially useful in e-commerce where small degradations can have immediate commercial consequences.

\section{Monitoring production ML systems}
Monitoring must cover more than CPU usage and uptime.
A live model can be operationally healthy and still commercially harmful.

\subsection{Technical monitoring}
Useful technical signals include:
\begin{itemize}
  \item latency,
  \item throughput,
  \item error rate,
  \item feature freshness,
  \item missing-value spikes,
  \item tool or dependency failures in inference paths.
\end{itemize}

\subsection{Model monitoring}
Useful model-focused signals include:
\begin{itemize}
  \item score distribution shifts,
  \item calibration changes,
  \item prediction drift,
  \item confidence anomalies,
  \item slice-level performance changes.
\end{itemize}

\subsection{Business monitoring}
Production evaluation should also include business outcomes such as:
\begin{itemize}
  \item click-through and conversion for recommenders,
  \item fraud-dollar capture and false-positive burden,
  \item margin and inventory effects for pricing models,
  \item containment and recontact rate for service assistants.
\end{itemize}

This is what turns monitoring from infrastructure tracking into operational accountability.

\section{Data drift and concept drift}
Drift is unavoidable in e-commerce ML.
Students should understand several forms of drift.

\subsection{Data drift}
Data drift occurs when the distribution of inputs changes.
Examples include:
\begin{itemize}
  \item seasonal traffic changes,
  \item new product categories,
  \item device mix shifts,
  \item changed promotion patterns,
  \item new fraud tactics.
\end{itemize}

\subsection{Concept drift}
Concept drift occurs when the relationship between inputs and outcomes changes.
For example:
\begin{itemize}
  \item the same browsing behavior stops indicating purchase intent after a UI redesign,
  \item fraud patterns evolve so past indicators weaken,
  \item pricing response changes because a competitor enters the market.
\end{itemize}

\subsection{Response to drift}
Possible responses include:
\begin{itemize}
  \item alerting,
  \item retraining,
  \item threshold adjustment,
  \item feature revision,
  \item temporary fallback to safer baselines.
\end{itemize}

\section{CI/CD for ML}
Continuous integration and deployment for ML differs from classical software delivery because the deployed artifact is influenced by data as well as code.

\subsection{What CI for ML should validate}
Reasonable validation steps include:
\begin{itemize}
  \item unit tests for transformations,
  \item schema and feature checks,
  \item reproducibility checks,
  \item training pipeline success,
  \item offline evaluation against baselines,
  \item safety or policy checks for sensitive outputs.
\end{itemize}

\subsection{What CD for ML should consider}
Deployment should consider:
\begin{itemize}
  \item compatibility with serving infrastructure,
  \item readiness of online features,
  \item staged rollout strategy,
  \item alert thresholds,
  \item rollback readiness.
\end{itemize}

In e-commerce, deployment speed matters, but so does release discipline.

\section{Experiment tracking and model comparison}
Teams need a reliable record of what was tried and what worked.

\subsection{What to track}
Useful tracked metadata include:
\begin{itemize}
  \item dataset reference,
  \item feature set version,
  \item model type,
  \item hyperparameters,
  \item training date and window,
  \item evaluation metrics,
  \item release decision.
\end{itemize}

\subsection{Why this matters}
Without structured experiment tracking, teams can neither compare models honestly nor explain why a deployed model was chosen.
This becomes particularly important when several business units or use cases share the same ML platform.

\section{Governance and approvals}
Not every model should be promoted automatically.
Some e-commerce use cases have a high enough impact that governance gates are necessary.

\subsection{Examples of governed models}
\begin{itemize}
  \item fraud and abuse scoring,
  \item dynamic pricing,
  \item customer-service agents with action tools,
  \item high-visibility recommendation systems,
  \item risk or compliance-related models.
\end{itemize}

\subsection{Typical governance artifacts}
Reasonable artifacts may include:
\begin{itemize}
  \item model cards,
  \item risk assessment notes,
  \item monitoring plans,
  \item evaluation summaries,
  \item rollback criteria,
  \item ownership and escalation paths.
\end{itemize}

Governance is not meant to slow all experimentation equally.
It is meant to scale decision quality for systems that matter.

\section{Incident response for ML systems}
Production ML systems need incident handling just like other production services.

\subsection{Examples of ML incidents}
\begin{itemize}
  \item a recommender starts surfacing unavailable items,
  \item a fraud model causes a spike in false declines,
  \item a service assistant cites the wrong policy version,
  \item a pricing model recommends below-cost prices after a cost feed failure.
\end{itemize}

\subsection{Preparedness}
Good preparedness includes:
\begin{itemize}
  \item clear runbooks,
  \item fallback modes,
  \item alert ownership,
  \item frozen-safe baselines,
  \item audit logs for recent model and data changes.
\end{itemize}

This is one reason operational simplicity is valuable.
Complex ML systems without response plans tend to fail expensively.

\section{Artifacts for this chapter}
The central artifact deliverable for this chapter is an ML system design package containing:
\begin{itemize}
  \item pipeline overview,
  \item feature and data dependencies,
  \item deployment path,
  \item monitoring plan,
  \item drift and rollback strategy,
  \item governance and ownership summary.
\end{itemize}

\section{Hands-on artifacts}
Students should produce:
\begin{itemize}
  \item an ML system design document,
  \item a monitoring plan with technical, model, and business signals,
  \item and a short incident playbook for one selected use case.
\end{itemize}

\section{Managerial implications}
Organizations often underestimate the operational cost of ML.
The model may be only a small fraction of the total system effort.
Most production pain comes from weak data contracts, poor observability, inconsistent features, and unclear ownership.

Well-run MLOps and DataOps allow the business to benefit from ML repeatedly rather than episodically.
They convert isolated model wins into an operational capability.

\section{Multiple-choice questions (MCQs)}
\begin{enumerate}
  \item What is \emph{training-serving skew}?
  \begin{enumerate}
    \item A difference between the fonts used in notebooks and dashboards.
    \item A mismatch between how features or transformations are produced in training and how they are produced at serving time.
    \item A pricing discount given only during model training.
    \item A mandatory property of all neural networks.
  \end{enumerate}

  \item Why is a model registry or equivalent release control useful?
  \begin{enumerate}
    \item It guarantees no model will ever fail.
    \item It helps track approved model versions, deployment state, evaluation evidence, and rollback options.
    \item It replaces monitoring entirely.
    \item It only matters for academic experiments.
  \end{enumerate}

  \item Which is the best example of a \emph{business} monitoring signal for an e-commerce ML system?
  \begin{enumerate}
    \item CPU temperature
    \item Fraud-dollar capture or revenue-per-session lift
    \item Container image size
    \item Disk inode count
  \end{enumerate}

  \item Data drift differs from concept drift because:
  \begin{enumerate}
    \item data drift is about input-distribution change, while concept drift is about change in the relationship between inputs and outcomes.
    \item concept drift is only about network outages.
    \item data drift only happens in recommendation systems.
    \item there is no real difference in practice.
  \end{enumerate}

  \item Why do production ML systems need incident playbooks?
  \begin{enumerate}
    \item Because model failures can affect real business outcomes and require fast, structured response.
    \item Because playbooks improve offline accuracy.
    \item Because incidents only happen in ERP systems.
    \item Because monitoring removes the need for rollback.
  \end{enumerate}
\end{enumerate}

\subsection*{Answer key}
\begin{enumerate}
  \item (b)
  \item (b)
  \item (b)
  \item (a)
  \item (a)
\end{enumerate}
